\begin{table}[!tbp]\centering
\caption{La crisis política de 2022--23 y el diseño de exposición regional continua}\label{tab:crisis}
\begin{threeparttable}
\footnotesize
\setlength{\tabcolsep}{3pt}
\begin{tabular}{lcccccc}
\toprule
& (1) Base & (2) +Post$\times$ & (3) +2023$\times$ & (4) Excl.\ bloque & (5) Post = & (6) (4)+(5) \\
Resultado & & protesta & protesta & sur & solo 2022 & \\
\midrule
Log salario/hora & 0.017 & 0.012 & 0.015 & 0.009 & 0.019 & 0.008 \\
 & (0.012) & (0.012) & (0.012) & (0.012) & (0.016) & (0.014) \\
 & [0.034] & [0.074] & [0.038] & [0.156] & [0.088] & [0.355] \\
Empleo & -0.034$^{***}$ & -0.027$^{***}$ & -0.029$^{***}$ & -0.029$^{**}$ & -0.029$^{***}$ & -0.024$^{**}$ \\
 & (0.012) & (0.009) & (0.009) & (0.010) & (0.009) & (0.009) \\
 & [0.001] & [0.004] & [0.001] & [0.000] & [0.002] & [0.007] \\
Empleo formal & -0.018$^{***}$ & -0.017$^{***}$ & -0.016$^{***}$ & -0.016$^{***}$ & -0.006 & -0.003 \\
 & (0.003) & (0.003) & (0.002) & (0.003) & (0.004) & (0.003) \\
 & [0.001] & [0.001] & [0.000] & [0.002] & [0.453] & [0.877] \\
Trab. independiente & 0.010$^{**}$ & 0.008$^{**}$ & 0.008$^{**}$ & 0.010$^{**}$ & 0.007$^{*}$ & 0.008 \\
 & (0.004) & (0.004) & (0.003) & (0.004) & (0.004) & (0.006) \\
 & [0.023] & [0.052] & [0.044] & [0.060] & [0.069] & [0.151] \\
\bottomrule
\end{tabular}
\begin{tablenotes}\footnotesize
\item Notas: Cada celda reporta el coeficiente de $\mathrm{Post}\times$(índice de Kaitz departamental estandarizado) del diseño de exposición continua, con errores estándar agrupados por departamento entre paréntesis y, entre corchetes, el valor $p$ de inferencia por aleatorización (permutando los índices de Kaitz entre los departamentos incluidos, 1000 extracciones, con el control de protestas fijo en su asignación real). La columna (1) es la especificación de base del Cuadro~\ref{tab:inference}, columna (5). Las columnas (2) y (3) añaden la movilización máxima de protesta per cápita estandarizada, interactuada con el indicador posterior a la reforma y con un indicador desde 2023Q1, respectivamente; la intensidad de las protestas es el número máximo de personas movilizadas en cada departamento en enero--febrero de 2023, según los reportes mensuales de conflictos de la Presidencia del Consejo de Ministros, dividido entre la población en edad de trabajar del departamento. La columna (4) excluye los seis departamentos del sur que fueron el centro de las protestas de diciembre de 2022 a marzo de 2023 (Apur\'imac, Arequipa, Ayacucho, Cusco, Madre de Dios, Puno). La columna (5) restringe la ventana posterior a la reforma a 2022Q3--Q4, antes de la destitución del presidente Castillo el 7 de diciembre de 2022, y la columna (6) combina ambas restricciones. Todas las estimaciones están ponderadas con los factores de expansión de la ENAHO. $^{*}p<0.1$, $^{**}p<0.05$, $^{***}p<0.01$.
\end{tablenotes}
\end{threeparttable}\end{table}
